% One-time imported draft. Editable native source; no Markdown regeneration.
\section{State the invariant at a precise
boundary}\label{005-seven-segment-counter-p07-state-the-invariant-at-a-precise-boundary}

At the start of an iteration, let j be the number of outputs already
completed. The index is f + j, the cursor is glyphs + f + j, and the
remaining count is n \ensuremath{-} j. The output history is exactly the first j
bytes of the requested window. For an entered iteration, the pre-test
guarantees that at least one remains.

The lookup and store append the next required byte. The cursor and index
then advance, and the remaining count decreases. At
\texttt{iteration\_done}, the same relationship holds with j increased
by one. When the count reaches zero, the requested history is complete.

The intermediate \texttt{displayed} label comes before those updates. It
shows the just-written glyph with the old index and count still present.
A debugger observation must name its boundary; otherwise two correct
snapshots can appear to disagree about the remaining count.
